% =====================================================================================
% Note de themes de travail, demandee au point du 13 aout 2026 (bilan de carriere).
% Destinataire : Hugo Rapior. A transmettre avant le point RH du 4 septembre, la liste
% devant ensuite etre communiquee a Axel et Ali.
%
% OBJECTIF DE CETTE VERSION : la note ne dit plus seulement ce que je veux faire, elle dit
% ce que le cabinet y gagne. Un bilan de carriere se decide sur la lisibilite du profil et
% sur l'actif qu'il laisse derriere lui, pas sur une liste de preferences.
%
% TROIS CONTRAINTES DE FORME, VOLONTAIRES :
%   - deux pages au maximum, et rien de decoratif ;
%   - AUCUN tiret cadratin. Le babel francais en met un en puce d'itemize par defaut, d'ou
%     le label redefini plus bas. A verifier sur le PDF, jamais sur la source ;
%   - aucun chiffre qui ne soit verifiable. Tous ceux de la section « preuves » se relevent
%     sur le depot du memoire.
%
% BLOC A COMPLETER PAR KELIAN : la section « Missions Nexialog », laissee en commentaire
% plus bas. C'est le seul endroit ou la note manque de matiere, et c'est le plus important
% aux yeux d'un manager. Le modele et la consigne y sont ecrits.
% =====================================================================================
\documentclass[11pt,a4paper]{article}

\usepackage[utf8]{inputenc}
\usepackage[T1]{fontenc}
\usepackage{lmodern}
\usepackage[french]{babel}
\usepackage[top=1.7cm,bottom=1.6cm,left=1.9cm,right=1.9cm]{geometry}
\usepackage{booktabs}
\usepackage{array}
\usepackage{enumitem}
\usepackage{xcolor}
\usepackage{titlesec}

% Charte Nexialog : les memes roles que style_nexialog.py, et aucun hexadecimal en clair
% ailleurs que dans ces quatre definitions.
\definecolor{encre}{HTML}{1B1E30}
\definecolor{accent}{HTML}{A6002E}
\definecolor{bleu}{HTML}{204993}
\definecolor{grisclair}{HTML}{F2F2F2}

\color{encre}
\titleformat{\section}{\normalfont\bfseries\color{bleu}}{}{0pt}{}
\titlespacing{\section}{0pt}{10pt}{3pt}
% Puce explicite : le babel francais mettrait un tiret cadratin.
\setlist[itemize]{label=\textcolor{accent}{\textbullet},itemsep=1.5pt,topsep=2pt,
                  leftmargin=1.1em,parsep=0pt}
\renewcommand{\arraystretch}{1.25}
\newcolumntype{L}[1]{>{\raggedright\arraybackslash}p{#1}}

\newcommand{\encadre}[2]{%
  \vspace{3pt}\noindent\colorbox{grisclair}{%
  \begin{minipage}{\dimexpr\textwidth-2\fboxsep}
  \footnotesize\textbf{#1} #2
  \end{minipage}}\vspace{3pt}}

\pagestyle{empty}

\begin{document}
\small

{\Large\bfseries\color{bleu} Thèmes de travail souhaités}\\[3pt]
{\color{accent}\rule{\textwidth}{1.1pt}}\\[3pt]
{\footnotesize Kélian Kaddouri \hfill Pour Hugo Rapior, avant le point RH du 4 septembre
2026}

\vspace{5pt}

\noindent Suite du bilan de carrière. La liste sépare ce que je veux \textbf{développer} de
ce que je suis prêt à \textbf{porter}, selon la distinction posée en séance. J'ai ajouté ce
que chaque thème ouvre pour le cabinet, parce que c'est le critère qui compte davantage que
ma préférence.

\vspace{3pt}
\noindent\textbf{En une phrase.} Devenir le référent quantitatif du cabinet sur le risque
cyber, en restant capable de tenir les sujets non-vie classiques, provisionnement et
capital, sans temps de montée en charge. C'est chez Nexialog que je veux construire cette
trajectoire.

\section{Ce que je veux développer}

\begin{center}\footnotesize
\begin{tabular}{@{}L{2.5cm} L{6.1cm} L{7.6cm}@{}}
\toprule
\textbf{Thème} & \textbf{Pourquoi celui-là} & \textbf{Ce que j'y apporte déjà} \\
\midrule
\textbf{R\&D cyber}
& Sujet jeune, donnée rare et mal partagée, aucune méthode de référence établie. C'est là
  qu'un travail de recherche a le plus de valeur marginale, et l'obligation est
  réglementaire depuis janvier 2025 sans qu'aucun module de formule standard n'existe.
& Mémoire de quantification d'un besoin de capital ORSA au titre de DORA par une cascade
  dirigée entre les cinq piliers : calibration valeurs extrêmes, identification partielle
  de la matrice de propagation, et un dispositif qui vérifie chaque nombre publié contre le
  script qui le produit. \\
\addlinespace
\textbf{Non-vie}
& La branche où vivent les outils que je maîtrise, fréquence et sévérité, queues lourdes,
  charge annuelle agrégée. Le cyber en est une composante, pas une exception : c'est ce qui
  rend le profil staffable en dehors de sa spécialité.
& Fréquence et sévérité modélisées séparément, tests d'adéquation, surdispersion, mesures
  de risque sur charge agrégée et allocation par Euler. \\
\addlinespace
\textbf{Réassurance et accumulation}
& Le prolongement direct de mon sujet : dès qu'on parle accumulation et cause commune, la
  question devient une question de structure de traité, par risque ou par événement. Et la
  cession est le premier levier qu'un client actionne avant le capital, y compris quand elle
  se traite dans un dossier non-vie sans mandat de réassurance dédié.
& Accumulation par cause commune, dépendance de queue, principe du grand saut unique, et le
  pont entre quantile d'un sinistre et quantile de la charge annuelle. \\
\bottomrule
\end{tabular}
\end{center}

\vspace{2pt}
\noindent\textbf{Ce que ces trois thèmes ouvrent pour le cabinet.}
\begin{itemize}
\item \textbf{Une offre outillée là où il n'y a pas de méthode de référence.} DORA
      s'applique depuis janvier 2025, aucun module ne le couvre en formule standard, et les
      assureurs doivent malgré tout le traiter en ORSA. La demande existe avant la méthode,
      ce qui est la situation la plus favorable pour un cabinet de conseil.
\item \textbf{Un actif réutilisable, pas un livrable unique.} Ce qui sort du mémoire n'est
      pas un chiffre mais une chaîne de calcul : un client la rejoue sur son propre bilan.
      Je peux la transformer en note méthodologique interne et en démonstration sur un
      bilan anonymisé, à un coût faible puisque le code existe.
\item \textbf{De la traçabilité, qui est un argument commercial et pas une coquetterie.}
      Chaque nombre est rattaché au script qui le calcule. Devant un régulateur ou un
      auditeur, c'est ce qui distingue un modèle défendable d'un tableur.
\item \textbf{De la visibilité.} Le rapport à venir de la direction générale du Trésor sur
      le cyber et les travaux de l'Institut des Actuaires sont deux points d'entrée, et des
      échanges sont déjà engagés de ce côté.
\end{itemize}

\section{Ce que je suis prêt à porter}

\noindent Sans en faire un axe, et sans réserve sur l'engagement.

\begin{itemize}
\item \textbf{Data science appliquée à la fraude.} Données déséquilibrées, arbitrage entre
      pouvoir prédictif et explicabilité. Le rapprochement avec le cyber est technique et
      non cosmétique : dans les deux cas l'événement est rare, la donnée est censurée par
      la détection, et ce qui n'est pas vu compte autant que ce qui est vu.
\item \textbf{Provisionnement.} Technique et normé, utile à connaître de l'intérieur pour
      ne pas rester cantonné au capital, et c'est le sujet qui donne accès aux directions
      techniques.
\item \textbf{Modélisation du capital.} Formule standard, ORSA, SCR par module. C'est déjà
      le terrain de mon mémoire, et c'est le langage commun avec les mêmes interlocuteurs.
\item \textbf{Gestion des risques.} Cartographie, criticité, appétit au risque, et
      l'articulation avec l'ORSA. C'est d'où vient mon sujet : la cascade entre les cinq
      piliers a été une cartographie qualitative avant d'être un modèle stochastique, et
      j'ai gardé les deux lectures. Le prolongement naturel est le dispositif et la fonction
      gestion des risques, pas seulement le chiffre qui en sort.
\item \textbf{IA générative.} Je l'utilise au quotidien pour produire et pour relire, et
      l'angle qui m'intéresse professionnellement est de l'encadrer comme un risque de
      modèle, avec la même exigence de traçabilité que pour un modèle de capital : savoir ce
      qui a été produit par quoi, et pouvoir le rejouer. Je ne le revendique pas comme une
      expertise, je le porte volontiers.
\end{itemize}

\encadre{Sur le point soulevé en séance.}{J'ai bien noté qu'une partie des sujets se
prennent sans passion, parce qu'ils financent la R\&D et légitiment le reste. Je n'ai pas de
difficulté avec cela et je ne cherche pas à choisir mes missions. Ce à quoi je tiens est
plus modeste : qu'une part identifiée de mon temps reste sur le cyber, pour que la
compétence continue de se construire au lieu de s'arrêter avec le mémoire. En échange, ce
temps produit quelque chose que le cabinet garde, et pas seulement un mémoire soutenu.}

\section{Ce sur quoi je suis staffable sans montée en charge}

\noindent Le risque d'un profil de recherche est de n'être plaçable que sur son sujet. Ce
qui suit est immédiatement mobilisable.

\begin{itemize}
\item \textbf{Solvabilité II} : SCR par module, besoin de capital en ORSA, allocation par
      entité, et la calibration des lois de perte qui va avec, sur données incomplètes.
\item \textbf{Industrialisation en Python} : une chaîne de plus de quatre-vingt-dix scripts
      versionnés, modules partagés, sorties reproductibles à l'octet entre deux postes.
      Habitude du versionnement et de la revue.
\item \textbf{Restitution} : rédaction technique en \LaTeX{}, une cinquantaine de figures
      produites sous charte graphique, et l'habitude de présenter un résultat toutes les
      semaines à un interlocuteur qui le remet en cause.
\end{itemize}

% =====================================================================================
% A COMPLETER, ET C'EST LE POINT LE PLUS IMPORTANT DE LA NOTE.
% Hugo a explicitement demande en seance que les descriptions disent le CONTEXTE et le
% GAIN, en citant l'exemple du « gain sur l'Ains Index ». Cette note est plus credible avec
% deux ou trois missions Nexialog qu'avec dix lignes sur le memoire.
% Modele a suivre pour chaque ligne, dans cet ordre :
%   client ou secteur (anonymise si besoin) ; ce qui etait demande ; ce que J'AI fait ;
%   le resultat CHIFFRE ou la decision que cela a permis de prendre.
% Decommenter le bloc ci-dessous et remplir. Retirer les crochets.
%
% \section{Ce que j'ai livre chez Nexialog}
% \begin{itemize}
% \item \textbf{[Secteur, type de mission].} [Ce qui etait demande.] [Ce que j'ai fait.]
%       [Resultat chiffre : gain, delai tenu, decision permise.]
% \item \textbf{[Secteur, type de mission].} [...]
% \item \textbf{[Secteur, type de mission].} [...]
% \end{itemize}
% =====================================================================================

\section{Les CV, et une question de format}

\noindent Les trois versions sont en préparation pour début septembre : le \textbf{complet}
(toutes expériences et projets académiques, autant de pages que nécessaire), qui sert de
référence, alimente les deux autres et part dans l'outil de Knowledge Management sur Dust ;
le \textbf{one-pager généraliste} ; et les \textbf{CV par thème}, un cyber, data et fraude,
un actuariat, chaque projet rangé là où il est le plus pertinent.

\vspace{2pt}
\noindent Deux points relevés en séance et pris en compte : les descriptions diront le
\textbf{contexte} et le gain obtenu, pas seulement la technique employée, et je demande à
Noemi un exemple du format PowerPoint attendu avant de mettre en forme, pour éviter de
refaire trois documents. Puisque les CV alimentent un outil qui extrait les compétences par
modèle de langage, je veillerai à ce que les mots-clés de chacun des thèmes ci-dessus y
figurent en clair plutôt que sous-entendus.

\vspace{6pt}
\noindent{\color{accent}\rule{\textwidth}{0.6pt}}\\[1pt]
{\footnotesize\itshape Liste à valider avant le point RH du 4 septembre, pour transmission
à Axelle et Ali.}

\end{document}
